\begin{frame}{자주 하는 실수 (2/2)}
  \begin{itemize}
    \item \kb{$\gamma = 1$ 을 아무 고민 없이 쓴다.} 축소 사상
      성질 $\|TV_1 - TV_2\| \le \gamma \|V_1 - V_2\|$ 이
      $\gamma = 1 < 1$ 이 \textbf{아니므로} 성립하지 않고
      수렴이 보장되지 않는다(유한 에피소드만). 보상이 무한히
      계속되는 문제에서 리턴 자체가 \textbf{발산}. ``할인 없이
      다 더하기''는 \textbf{보상 합이 유한}한 문제에만 안전.
    \item \kb{``sweep 1회 = 하나의 상태만 갱신''이라고 오해.}
      sweep = \textbf{모든 상태 한 바퀴}(\texttt{for s in
      range(n\_states)} 한 바퀴). ``sweep 100회'' =
      $n\_states \times 100$ 번의 \textbf{개별 갱신}. sweep 수와
      개별 갱신 수를 혼동하면 계산량($n\_states$ 배)을
      \textbf{과소평가}.
  \end{itemize}
\end{frame}
